\documentclass[11pt]{article}
\usepackage[margin=1in]{geometry}
\usepackage{amsmath, amssymb, amsthm, mathtools}
\usepackage{hyperref}

\newtheorem{theorem}{Theorem}
\newtheorem{lemma}[theorem]{Lemma}
\newtheorem{corollary}[theorem]{Corollary}
\theoremstyle{definition}
\newtheorem{definition}{Definition}
\theoremstyle{remark}
\newtheorem*{remark}{Remark}

\DeclareMathOperator{\Var}{Var}
\DeclareMathOperator{\Cov}{Cov}
\DeclareMathOperator{\Cor}{Cor}
\newcommand{\E}{\mathbb{E}}
\newcommand{\R}{\mathbb{R}}

\title{Variance and Covariance}
\author{math-foundations / 26-variance-covariance}
\date{}

\begin{document}
\maketitle

\section{Setup}

Throughout this lesson, fix a probability space $(\Omega, \mathcal{F}, P)$ and let
$L^2(P) = \{X : \Omega \to \R \mid \E[X^2] < \infty\}$ denote the Hilbert space of
square-integrable random variables. All random variables below are assumed to
lie in $L^2(P)$, which guarantees the existence of the moments $\E[X]$, $\E[X^2]$,
and (by Cauchy--Schwarz) $\E[XY]$ for $X, Y \in L^2(P)$.

\section{Definitions}

\subsection*{First, an example}
Before the formal definitions, consider $X \sim \mathrm{Bernoulli}(p)$ with $p = 0.3$:
$X$ takes value $1$ with probability $0.3$ and $0$ with probability $0.7$. Then
$\E[X] = 0.3$ and $\E[X^2] = 0.3$, so
\[
\Var(X) = \E[X^2] - (\E[X])^2 = 0.3 - 0.09 = 0.21 = p(1 - p).
\]
\emph{Read aloud:} the variance of X equals expectation of X squared minus the square of the expectation of X, which is zero point three minus zero point zero nine, equals zero point two one, equals p times one minus p.
Now let $Y = 2X$. Intuitively $Y$ moves in lockstep with $X$, just twice as far, so they
should be perfectly positively correlated. Indeed
\[
\Cov(X, Y) = \Cov(X, 2X) = 2\,\Var(X) = 0.42, \qquad \rho_{X,Y} = \frac{0.42}{\sqrt{0.21}\,\sqrt{4 \cdot 0.21}} = 1.
\]
\emph{Read aloud:} covariance of X and Y equals covariance of X with two X equals two times variance of X equals zero point four two; and the correlation rho sub X Y equals zero point four two divided by square root of zero point two one times square root of four times zero point two one, which equals one.
The rest of this lesson develops the definitions and identities behind these computations.

\begin{definition}[Variance]
The \emph{variance} of $X \in L^2(P)$ is
\[
\Var(X) \;\coloneqq\; \E\!\left[(X - \E[X])^2\right].
\]
\emph{Read aloud:} the variance of X is defined as the expectation of the quantity X minus the expectation of X, all squared.
By expanding the square and applying linearity of expectation,
\[
\Var(X) = \E[X^2] - (\E[X])^2.
\]
\emph{Read aloud:} the variance of X equals the expectation of X squared minus the square of the expectation of X.
The \emph{standard deviation} is $\sigma_X = \sqrt{\Var(X)}$.
\end{definition}

\begin{definition}[Covariance]
The \emph{covariance} of $X, Y \in L^2(P)$ is
\[
\Cov(X, Y) \;\coloneqq\; \E\!\left[(X - \E[X])(Y - \E[Y])\right]
           \;=\; \E[XY] - \E[X]\,\E[Y].
\]
\emph{Read aloud:} the covariance of X and Y is defined as the expectation of the product of X minus the expectation of X times Y minus the expectation of Y, which equals the expectation of X times Y minus the expectation of X times the expectation of Y.
Note $\Cov(X, X) = \Var(X)$.
\end{definition}

\begin{definition}[Correlation]
If $\sigma_X, \sigma_Y > 0$, the \emph{(Pearson) correlation coefficient} of $X$ and $Y$ is
\[
\rho_{X, Y} \;\coloneqq\; \frac{\Cov(X, Y)}{\sigma_X \, \sigma_Y}.
\]
\emph{Read aloud:} rho sub X comma Y is defined as the covariance of X and Y divided by the product of sigma sub X and sigma sub Y.
By Cauchy--Schwarz applied to the centered variables, $\rho_{X, Y} \in [-1, 1]$.
\end{definition}

\section{Elementary properties}

\begin{lemma}[Bilinearity and symmetry of covariance]
For $a, b, c \in \R$ and $X, Y, Z \in L^2(P)$:
\begin{enumerate}
\item $\Cov(X, Y) = \Cov(Y, X)$.
\item $\Cov(aX + b, Y) = a\,\Cov(X, Y)$.
\item $\Cov(X + Z, Y) = \Cov(X, Y) + \Cov(Z, Y)$.
\end{enumerate}
\end{lemma}

\begin{corollary}[Affine scaling of variance]
For $a, b \in \R$, $\Var(aX + b) = a^2 \, \Var(X)$.
\end{corollary}

\section{Variance of a sum (main theorem)}

\begin{theorem}[Variance of a sum]
\label{thm:var-sum}
For $X, Y \in L^2(P)$,
\[
\Var(X + Y) \;=\; \Var(X) \;+\; \Var(Y) \;+\; 2\,\Cov(X, Y).
\]
\emph{Read aloud:} the variance of X plus Y equals the variance of X plus the variance of Y plus two times the covariance of X and Y.
\end{theorem}

\begin{proof}
Let $\mu_X = \E[X]$ and $\mu_Y = \E[Y]$. By linearity of expectation,
$\E[X + Y] = \mu_X + \mu_Y$. Applying the definition of variance to $X + Y$,
\begin{align*}
\Var(X + Y)
  &= \E\!\left[\,\big((X + Y) - \E[X + Y]\big)^2\,\right] \\
  &= \E\!\left[\,\big((X - \mu_X) + (Y - \mu_Y)\big)^2\,\right] \\
  &= \E\!\left[\,(X - \mu_X)^2 + 2(X - \mu_X)(Y - \mu_Y) + (Y - \mu_Y)^2\,\right].
\end{align*}
\emph{Read aloud:} the variance of X plus Y equals the expectation of the square of the quantity X plus Y minus the expectation of X plus Y, which equals the expectation of the square of the centered sum X minus mu sub X plus Y minus mu sub Y, which expands to the expectation of X minus mu sub X squared plus twice X minus mu sub X times Y minus mu sub Y plus Y minus mu sub Y squared.
By linearity of expectation,
\begin{align*}
\Var(X + Y)
  &= \E\!\left[(X - \mu_X)^2\right]
   + 2\,\E\!\left[(X - \mu_X)(Y - \mu_Y)\right]
   + \E\!\left[(Y - \mu_Y)^2\right] \\
  &= \Var(X) + 2\,\Cov(X, Y) + \Var(Y),
\end{align*}
\emph{Read aloud:} the variance of X plus Y equals the expectation of X minus mu sub X squared plus two times the expectation of X minus mu sub X times Y minus mu sub Y plus the expectation of Y minus mu sub Y squared, which equals the variance of X plus two times the covariance of X and Y plus the variance of Y.
where the last equality uses the definitions of variance and covariance.
\end{proof}

\begin{corollary}[$n$-term version]
For $X_1, \dots, X_n \in L^2(P)$,
\[
\Var\!\left(\sum_{i=1}^n X_i\right)
  = \sum_{i=1}^n \Var(X_i) + 2\!\!\sum_{1 \le i < j \le n}\!\! \Cov(X_i, X_j).
\]
\emph{Read aloud:} the variance of the sum from i equals one to n of X sub i equals the sum from i equals one to n of the variance of X sub i, plus two times the sum over pairs i less than j from one to n of the covariance of X sub i and X sub j.
In particular, if the $X_i$ are pairwise uncorrelated, then
$\Var(\sum X_i) = \sum \Var(X_i)$.
\end{corollary}

\begin{remark}
Independence of $X$ and $Y$ implies $\E[XY] = \E[X]\E[Y]$, hence
$\Cov(X, Y) = 0$. The converse is false: take $X \sim \text{Uniform}(\{-1, 0, 1\})$
and $Y = X^2$; then $\Cov(X, Y) = 0$ but $Y$ is a deterministic function of $X$.
\end{remark}

\section{Worked example}

\paragraph{Bernoulli.} Let $X \sim \mathrm{Bernoulli}(p)$. Then $\E[X] = p$,
$\E[X^2] = p$, so
\[
\Var(X) = p - p^2 = p(1 - p).
\]
\emph{Read aloud:} the variance of X equals p minus p squared, which equals p times one minus p.
The variance is maximized at $p = 1/2$.

\end{document}
